\chapter{Tag 2 — ISBN-10, Luhn, Hamming-Distanz}

Stub-Kapitel für M1. Inhalte folgen in M4.

Dieses Kapitel testet \emph{Szenario 4}: ein Recto-/Verso-Wechsel mitten
im Kapitel — die zweite Aufgabe muss automatisch die Spiegel-Variante
bekommen.

\begin{bleistiftuebung}{1}{Modulo warmrechnen}
Berechne im Kopf oder auf Papier: \(29 \bmod 7\), \(100 \bmod 9\),
\(1000 \bmod 11\), \((47 + 38) \bmod 10\).

Modulo \(m\) heißt: teile durch \(m\) und nimm den \emph{Rest}. Du wirst
es heute oft brauchen.
\end{bleistiftuebung}

Etwas Fließtext, damit die nächste Seite einen klaren Anfang hat.

\clearpage

\begin{pythoneinheit}{1}{EAN-13-Validator}
Implementiere die Funktion zur Berechnung der EAN-13-Prüfziffer:

\begin{verbatim}
def ean13_pruefziffer(zwoelf_ziffern):
    summe = sum(d * (1 if i % 2 == 0 else 3)
                for i, d in enumerate(zwoelf_ziffern))
    return (10 - summe % 10) % 10
\end{verbatim}

Probiere die Funktion mit der EAN auf einer Schokoladenpackung aus.
\end{pythoneinheit}
